% BHU PYQ -- Professional Exam Template
\documentclass[11pt,a4paper]{article}

\usepackage[a4paper,top=2.5cm,bottom=2.5cm,left=2.8cm,right=2.8cm,headheight=14pt]{geometry}
\usepackage{amsmath,amssymb}
\usepackage[T1]{fontenc}
\usepackage[utf8]{inputenc}
\usepackage{lmodern}
\usepackage{microtype}
\usepackage{fancyhdr}
\usepackage{enumitem}
\usepackage{listings}
\usepackage{hyperref}
\usepackage{lastpage}
\usepackage{parskip}

\hypersetup{colorlinks=true,urlcolor=black,linkcolor=black,
  pdftitle={CSB-501: Object Oriented Programming Using C++ -- BHU PYQ}}

\pagestyle{fancy}
\fancyhf{}
\renewcommand{\headrulewidth}{0.4pt}
\renewcommand{\footrulewidth}{0.4pt}
\fancyhead[L]{\small   \textit{Banaras Hindu University}}
\fancyhead[R]{\small   \textit{CSB-501: Object Oriented Programming Us}}
\fancyfoot[C]{\small Page  \thepage\ of \pageref{LastPage}}

\lstset{
  basicstyle=\small tfamily,
  frame=single,
  breaklines=true,
  columns=flexible,
  keepspaces=true
}


\newlist{parts}{enumerate}{2}
\setlist[parts,1]{label=(\alph*),leftmargin=*,itemsep=4pt,topsep=3pt}
\setlist[parts,2]{label=(\roman*),leftmargin=*,itemsep=2pt,topsep=2pt}

\newcommand{\pts}[1]{\hfill   \textit{[#1]}}


\begin{document}

\begin{center}
  {\large\bfseries BANARAS HINDU UNIVERSITY}\\[2pt]
  {\normalsize B.Sc. (Hons.) Semester V Examination, 2013-14}\\[6pt]
  \rule{0.8\linewidth}{0.5pt}\\[6pt]
  {\large\bfseries Paper: CSB-501 --- Object Oriented Programming Using C++}\\[4pt]
  {\normalsize Time: Three Hours \hfill Maximum Marks: 70}
\end{center}

\rule{\linewidth}{0.4pt}

\smallskip



\noindent   \textit{Instructions: Answer any FIVE questions including Question 1 which is compulsory. The figures in the right-hand margin indicate full marks.}

\bigskip

\medskip


\noindent   \textbf{Question 1 (Compulsory).}\pts{5 + 5 + 4 = 14}

\begin{parts}
  \item State whether the following statements are TRUE or FALSE:
  
\begin{enumerate}[label=(\roman*),itemsep=2pt]
    \item Polymorphism is extensively used in implementing inheritance.
    \item In C++, a class is a user-defined data type.
    \item The scope resolution operator can be used to uncover a hidden variable.
    \item The default visibility mode in C++ is private.
    \item In C++, both virtual constructors and destructors are allowed.
  \end{enumerate}
  
  \item Fill in the blanks:
  
\begin{enumerate}[label=(\roman*),itemsep=2pt]
    \item The operator   \texttt{<<} is called the \_\_\_\_\_\_ operator.
    \item A \_\_\_\_\_\_ takes a reference to an object of the same class as itself as an argument.
    \item An \_\_\_\_\_\_ class is one that is not used to create objects.
    \item C++ uses a unique keyword called \_\_\_\_\_\_ to represent an object that invokes a member function.
    \item C++ supports a mechanism known as \_\_\_\_\_\_ to achieve runtime polymorphism.
  \end{enumerate}
  
  \item Find errors, if any, in the following C++ statements:
  
\begin{enumerate}[label=(\roman*),itemsep=2pt]
    \item   \texttt{int *p = new;}
    \item   \texttt{enum \{green, yellow, red\};}
    \item   \texttt{for (i=1; int i<10; i++) cout << i << " extbackslash{}n";}
    \item   \texttt{char name[3] = "USA";}
  \end{enumerate}
\end{parts}

\bigskip


\noindent   \textbf{Question 2.}\pts{5 + 5 + 4 = 14}

\begin{parts}
  \item What are the unique advantages of an object-oriented programming paradigm?
  \item Write the output of the following C++ program:

\begin{lstlisting}[language=C++]
#include <iostream>
using namespace std;
int a = 20;
int main() {
    int a = 40;
    {
        int k = a;
        int a = 50;
        cout << "INNER block
";
        cout << "k=" << k << "
";
        cout << "a=" << a << "
";
        cout << "::a=" << ::a << "
";
    }
    cout << "OUTER BLOCK
";
    cout << "a=" << a << "
";
    cout << "::a=" << ::a << "
";
    return 0;
}
\end{lstlisting}
  \item What is a reference variable? What is the major use of this variable?
\end{parts}

\bigskip


\noindent   \textbf{Question 3.}\pts{6 + 3 + 5 = 14}

\begin{parts}
  \item What do you mean by overloading of a function? When do we use this concept? Explain with an example.
  \item When do we declare a member of a class static?
  \item Write a program in C++ that shows how to use a common friend function to exchange the private values of two classes.
\end{parts}

\bigskip


\noindent   \textbf{Question 4.}\pts{5 + 5 + 4 = 14}

\begin{parts}
  \item What is a constructor and a destructor? Also write the various types of constructors available in C++.
  \item Write any five rules for overloading operators.
  \item How many arguments are required in the definition of an overloaded unary operator? Explain with an example.
\end{parts}

\bigskip


\noindent   \textbf{Question 5.}\pts{5 + 5 + 4 = 14}

\begin{parts}
  \item Explain the effect of inheritance on the visibility of class members using a diagram.
  \item Write a program in C++ that demonstrates multilevel inheritance.
  \item Write any four rules for virtual functions in C++.
\end{parts}

\bigskip


\noindent   \textbf{Question 6.}\pts{4 + 5 + 5 = 14}

\begin{parts}
  \item Both   \texttt{cin} and   \texttt{getline()} functions can be used for reading a string. Comment.
  \item What is a file mode? Describe the various file mode options available in C++.
  \item Write a program in C++ that reads two files simultaneously.
\end{parts}

\bigskip


\noindent   \textbf{Question 7.} Write short notes on any two of the following:\pts{7 $ \times$ 2 = 14}

\begin{parts}
  \item Object-Oriented Analysis and Design (OOAD)
  \item Class Templates
  \item Exception Handling
  \item Class Wizard
\end{parts}

\vfill
\rule{\linewidth}{0.4pt}

\begin{center}\small
\href{https://bhu-exam.vercel.app/}{Click here to visit our website}\\[4pt]\href{https://me-aryan.vercel.app/}{Click here to visit our contributor}\\[4pt]\href{https://quashan-me.vercel.app/}{Click here to visit our contributor}
\end{center}
\end{document}